% Second pass (Claude, 2026-11-10): first pass checked against the archive images; changes in LM_C4-039-057_changes.md.
\scanpage{39}
\mspage{667}

\[
\tag{24} l_1^\nu=\lambda_1^\nu\omega^{\nu_2}.
\]

Posant (cf (19))
\[
\tag{25}
\left\{\begin{aligned}
\beta&=\beta_0\tau^{\nu_2}\qquad(\unsure{\beta_0|\beta}\in\Gamma\underline M)\\
\eta_0&=\beta_0\sigma(\beta_0)^{-1}
\end{aligned}\right.
\]
on aura donc
\[
\eta\overset{\unsure{\mathrm{dfn}}}{=}
\beta\sigma(\beta)^{-1}
=\beta_0\sigma(\beta_0)\omega^{\nu_2}
=\eta_0\omega^{\nu_2}
\]
et l'équation du lacet prend la forme
d'une équation dans $\underline M$
\[
\tag{26}
\unsure{\zeta|\xi}=\deleted{\unsure{\beta}}\eta_0\lambda_1^\nu,
\qquad\text{i.e. \unsure{membre} de (21)}
\]
\[
\tag{26bis}\unsure{\zeta|\xi}=
\]
où $\unsure{\zeta|\xi},\eta_0$ sont deux sections de $\underline M$ qui
se déterminent mutuellement (pour $\nu$
donné). On a ainsi exprimé de façon
commode l'équation du lacet, i.e. la
condition c)\inserted{de (17)}, il reste à exprimer a) et b),
qui se formulent aussi comme
\[
\tag{27}
\left\{\begin{aligned}
\text{a)}\;&\unsure{\zeta|\xi}\ \text{localement de la forme}\quad
\overset{\inserted{(\nu_3=0)}}{\alpha_0\rho(\alpha_0)^{-1}}
\quad\text{ou}\quad
\overset{\inserted{(\nu_3=1)}}{\alpha_0\lambda_1^{-1}\rho(\alpha_0)^{-1}}\\
&\text{avec }\alpha_0\in\Gamma\underline M\\
\text{b)}\;&\eta_0\ \text{localement de la forme}\quad\beta_0\sigma(\beta_0)^{-1}
\end{aligned}\right.
\]
(Conditions dont la formulation est indépendante
de $\nu,\nu_2$ \deleted{\unsure{pour|$\nu_3$})}, mais dépend \unsure{apparemment} de
$\nu_3$.) Pour expliciter ces conditions, on
va écrire $\underline M$ additivement, les relations
s'écrivent
\margin{On a donc posé aussi $\alpha=\alpha_0\tau^{\nu_3}$}
\[
\tag{28}
\left\{\begin{aligned}
\unsure{\zeta|\xi}&=\alpha_0-\rho(\alpha_0)
\quad\overset{\inserted{\text{cas}}}{(\nu_3=0)}
\quad\text{ou}\quad
\unsure{\zeta|\xi}=\alpha_0-\rho(\alpha_0)-\lambda_1
\quad\overset{\inserted{\text{cas}}}{(\nu_3=1)}\\
\eta_0&=\beta_0-\sigma(\beta_0),
\end{aligned}\right.
\]
\margin{38}
\margin{38\\\texttt{https://grothendieck.umontpellier.fr}\\Université de Montpellier, UMR CNRS 5149, Case courrier 051, Place Eugène Bataillon\\34095 Montpellier cedex 5 - France}

\scanpage{40}

ce qu'on exprime aussi par
\[
\tag{29}
\left\{\begin{aligned}
\unsure{\zeta|\xi}&\in\operatorname{Im}(\mathrm{Id}_{\underline M}-\rho_{\underline M})
\quad(\nu_3=0)\quad\text{ou}\quad
\unsure{\zeta|\xi}\in-\lambda_1+\operatorname{Im}(\mathrm{Id}_{\underline M}-\rho_{\underline M})\\
\eta_0&\in\operatorname{Im}(\mathrm{Id}_{\underline M}-\sigma_{\underline M}).
\end{aligned}\right.
\]

C'est le moment de passer à des calculs
matriciels explicites, en rapportant $\underline M$ à
la base $\lambda_0,\lambda_1$. Dans cette base, les
\unsure{opérations} de $\rho_{\underline M},\sigma_{\underline M},\tau_{\underline M}$ sont données par
\[
\tag{30}
\left\{\begin{aligned}
\rho_{\underline M}&=\begin{pmatrix}0&-1\\1&-1\end{pmatrix}
&\mathrm{id}_{\underline M}-\rho_{\underline M}&=\begin{pmatrix}1&1\\-1&2\end{pmatrix}\\
\sigma_{\underline M}&=\begin{pmatrix}0&1\\1&0\end{pmatrix}
&\mathrm{id}_{\underline M}-\sigma_{\underline M}&=\begin{pmatrix}1&-1\\-1&1\end{pmatrix}\\
\tau_{\underline M}&=\begin{pmatrix}-1&0\\0&-1\end{pmatrix}.&&
\end{aligned}\right.
\]

On a
\[
\det(\mathrm{id}_{\underline M}-\rho_{\underline M})=3,
\qquad\det(\mathrm{id}_{\underline M}-\sigma_{\underline M})=0
\]
ce qui montre que \deleted{$\mathrm{id}_{M}$} les endomorphismes
$\mathrm{id}_{\underline M}-\rho_{\underline M}$ et $\mathrm{id}_{\underline M}-\sigma_{\underline M}$ de $\underline M$ ne sont
pas tout à fait des automorphismes\inserted{i.e. des épimorphismes} ---
Ça déconne en car. 3 pour le premier, \uline{partout}
\deleted{\unsure{car. 2}}\inserted{(en tte caractéristique)} pour le second.
\deleted{On prendra donc}\inserted{On aura}
\deleted{plus} intérêt à prendre comme "paramètres"
non $\unsure{\zeta|\xi}$ et $\eta_0$, mais bien évidemment
$\alpha_0$ et $\beta_0$ (en plus des "paramètres" discrets
$\nu_2,\nu_3$ et $\nu'$), et de \deleted{définir}\inserted{faire} \deleted{donc} un
groupe (si possible\ldots) avec les systèmes

\margin{39}
\margin{39\\\texttt{https://grothendieck.umontpellier.fr}\\Université de Montpellier, UMR CNRS 5149, Case courrier 051, Place Eugène Bataillon\\34095 Montpellier cedex 5 - France}

\scanpage{41}
\mspage{668}

\deleted{\unsure{(24)}}
\[
\tag{31}
\widetilde{\underline Z}_{\rho,\sigma}
=\left\{(\nu,\nu',\nu_2,\nu_3,\alpha_0,\beta_0)\ \left|\
\begin{gathered}
\nu,\nu'\in\Gamma\mathbf G_a,\quad
\nu_2,\nu_3\in\Gamma\{0,1\}_{S_0},\quad
\alpha_0,\beta_0\in\Gamma\underline M\\
\left\{\begin{aligned}
\unsure{\nu}&=2\nu'\unsure{-|+}\nu_2\qquad 2\nu+1\in\Gamma\mathbf G_m\\
\underbrace{\alpha_0-\rho(\alpha_0)\unsure{-|+}\nu_3\lambda_1}_{\unsure{\zeta|\xi}}
&=\underbrace{\beta_0-\sigma(\beta_0)}_{\eta}+\nu\lambda_1
\end{aligned}\right.
\end{gathered}\right.\right\}
\]
\margin{\uline{NB} $\nu'$ détermine $\nu$ qui est redondant}

L'équation des lacets, en $\alpha_0,\beta_0$, prend donc
la forme
\[
\tag{32}
(\unsure{\mathrm{id}|1}-\rho_{\underline M})\alpha_0
\unsure{-|=}(\mathrm{id}-\sigma_{\underline M})(\beta_0)
=\bigl(\underbrace{2\nu'\unsure{-|+}\nu_2}_{\nu}\unsure{+|-}\nu_3\bigr)\lambda_1
\]

\margin{Cela doit en effet \unsure{provenir du} fait que l'hom.
\diagram{\begin{tikzcd}[column sep=large]
\underline M\times\underline M\arrow[r]&\underline M\\
(\alpha_0,\beta_0)\arrow[r,mapsto]&
(\mathrm{id}-\rho_{\underline M})\alpha_0-
(\deleted{\unsure{\mathrm{id}}}\inserted{1}-\sigma_{\underline M})(\beta_0)
\end{tikzcd}}
\unsure{soit|est} bien épi\ldots\ Mais ce n'est le cas qu'en car. $\ne3$ !}

\uline{Remarque} Le schéma $\underline Z_{\rho,\sigma}$ lui $\overline{\mathrm m}$, tel qu'il
était défini au \S\ 49 \deleted{\unsure{\S\ 6}} $\overline{\underline{\mathrm{VI}}}$, n'est pas lisse
il me semble --- \unsure{ça coince|se corse} en car. 2 et 3.
Par contre le groupe $\widetilde{\underline Z}_{\rho,\sigma}$ plus fin,
défini en mettant dans le coup les
$\alpha,\beta$ tout de suite (au lieu de $(\unsure{\zeta|\xi},\eta)$)
semble bien \uline{lisse}. \unsure{Ça} amène à nouveau
à modifier un peu la présentation des
constructions du \S\ 49, en \unsure{m'interdisant|n'utilisant}
\unsure{encore un} groupe $\underline Z_{\rho,\sigma}$ (lié aux coordonnées
$\unsure{\zeta|\xi},\eta$), et en le remplaçant par
un groupe qui \unsure{apparaîtra|apparaît} \inserted{\unsure{à la fois|de façon}} comme un
normalisateur de $L_0$ (ou de $L_0^{\unsure{\sharp}}$) dans $\mathcal M_{\rho,\sigma}$
(ou $\mathcal M_{\rho,\sigma}^{\unsure{\natural|\sharp}}$). On fera une présentation dans
ce sens dans le \deleted{\unsure{suite}} \S\ suivant, et on
\unsure{suppose} ici déjà acquise cette construction,



\margin{40}
\margin{40\\\texttt{https://grothendieck.umontpellier.fr}\\Université de Montpellier, UMR CNRS 5149, Case courrier 051, Place Eugène Bataillon\\34095 Montpellier cedex 5 - France}

\scanpage{42}

D'où le \deleted{\unsure{sphère|schéma}} (schéma \deleted{et) groupe}
$\widetilde{\underline Z}_{\rho,\sigma}$ aura donc
une structure de groupe convenable, et
opère sur $\unsure{\underline C^{+}|\underline G^{+}}$ via
\[
\tag{33}
\left\{\begin{aligned}
u(\rho)&=(\alpha_0-\rho(\alpha_0)+\nu_3\lambda_1)\rho\\
u(\sigma)&=(\beta_0-\sigma(\beta_0))\omega^{\nu_2}\sigma.
\end{aligned}\right.
\]

Mais il faut quand $\overline{\mathrm m}$ vérifier que
ces formules (33) déterminent bel et
bien un automorphisme $u$ de $\unsure{\underline C^{+}|\underline G^{+}}$. La
vérification essentielle : faire en sorte
que la \deleted{formule \unsure{[7]}}\inserted{relation} (7)
\[
\tag{34}[\sigma,\lambda_\infty]
\]
est respectée par $u$.

Je vais expliciter l'action de $u$ sur
$\underline M$, en explicitant ses valeurs sur
\[
\left\{\begin{aligned}
\lambda_0&=\omega l_0=\omega\sigma\rho\sigma\rho\\
\lambda_1&=\omega l_1=\omega\rho\sigma\rho\sigma
\end{aligned}\right.
\]
en termes de \deleted{$\unsure{\zeta|\xi}$ et $\eta_0$, \unsure{intro}}
\[
\tag{35}
\left\{\begin{aligned}
\unsure{\zeta|\xi}&=\alpha_0-\rho(\alpha_0)+\nu_3\lambda_1\\
\eta_0&=\beta_0-\sigma(\beta_0)
\end{aligned}\right.
\qquad\unsure{\alpha|\&}\qquad\in\Gamma\underline M
\]
Liés à $u$ par
\[
\tag{36}
u\rho=\unsure{\zeta|\xi}\rho,\qquad
u\sigma=\omega^{\nu_2}\eta_0\sigma.
\]
On aura \deleted{donc} par la relation des lacets
\[
u(\varepsilon_0)=\deleted{\unsure{\mu}}\varepsilon_0^{\deleted{\unsure{\mu}}\inserted{(\mu)}}
\qquad(\mu=2\nu+1=4\nu'+2\nu_2+1)
\]
d'où

\margin{41}
\margin{41\\\texttt{https://grothendieck.umontpellier.fr}\\Université de Montpellier, UMR CNRS 5149, Case courrier 051, Place Eugène Bataillon\\34095 Montpellier cedex 5 - France}

\scanpage{43}
\mspage{669}
\[
u(l_0)=u(\varepsilon_0^2)=l_0^{(\mu)}
\]
et
\[
u(\lambda_0)=\unsure{u}(\omega l_0)
=\underbrace{u(\omega)}_{=\omega}\,u(l_0)
=\omega l_0^{(\mu)}
=\underbrace{\omega l_0}_{\lambda_0}
 \bigl(\underbrace{l_0^{(2\nu)}}_{\lambda_0^{2\nu}}\bigr)
=\lambda_0^{(2\nu+1)}=\lambda_0^{(\mu)}
\]
comme il convient. (Pour me rassurer, j'ai vérifié que la deuxième formule (33) donne bien
\[
u(\sigma^2)=(\eta_0\sigma)(\eta_0\sigma)
=\bigl(\underbrace{\eta_0\,\sigma(\eta_0)}_{1}\bigr)
 \deleted{\unsure{\sigma^2}}\omega=\omega
\quad\text{i.e.}\quad\deleted{\text{et}}
\]
\[
\tag{37} u(\omega)=\omega
\]
On aura
\[
u(\lambda_1)=u(\rho(\lambda_0))
=(\unsure{\zeta|\xi}\rho u)(\lambda_0)
=(\unsure{\zeta|\xi}\rho)u(\lambda_0)
=(\unsure{\zeta|\xi}\rho)(\lambda_0^\mu)=\lambda_1^\mu
\]
donc finalement, comme prédit
\[
\tag{38} u(\lambda)=\mu.\lambda\qquad\text{pour}\quad\lambda\in\Gamma\underline M
\]
\margin{où $\mu=2\nu+1=4\nu'+2\nu_2+1$}
On aura bien que $u(\sigma)=(\beta_0-\sigma(\beta_0))\omega^{\nu_2}.\sigma$
commute avec $u(\lambda_\infty)=\mu.\lambda_\infty$.

Tous ces calculs sont un peu \unsure{vaseux}, logiquement, car à la fois je veux prouver
l'existence d'un \deleted{\unsure{hom}} $u$ \inserted{de $\unsure{\underline C^{+}|\underline G^{+}}$},
et dans les calculs précédents je l'admets déjà (\unsure{implicitement}, dans un groupe
``plus gros'', savoir le schéma en groupes \unsure{lisse|libre} \inserted{(\unsure{heuristique})}
engendré par $\rho,\sigma\ldots$.
Une façon de les rendre \uline{rigoureux} serait d'utiliser la description intrinsèque de
$\unsure{\underline C^{+}|\underline G^{+}}$ :
\deleted{de \unsure{démontrer}}
\[
\tag{39}
\unsure{\underline C^{+}|\underline G^{+}}
=\unsure{C^{+}_{S_0}|G^{+}_{S_0}}
 \mathop{\wedge}_{M_{S_0}}\underline M_{S_0}
\]
donc définir un hom\inserted{$u$} de
$\unsure{\underline C^{+}|\underline G^{+}}$ dans un $S_0$-groupe $G$
(ici $G=\unsure{\underline C^{+}|\underline G^{+}}$) revient à
\margin{42}
\margin{42\\\texttt{https://grothendieck.umontpellier.fr}\\Université de Montpellier, UMR CNRS 5149, Case courrier 051, Place Eugène Bataillon\\34095 Montpellier cedex 5 - France}

\scanpage{44}
définir des homomorphismes
\[
\tag{40}
\diagram{\begin{tikzcd}[column sep=large]
\underline M_{S_0}\arrow[r,"u_{\underline M}"] & \unsure{G^{\sharp}|G^{\natural}}
\end{tikzcd}}
,\quad
\diagram{\begin{tikzcd}[column sep=large]
\unsure{C_{S_0}|G_{S_0}}\arrow[r,"u_{\unsure{C|G}}"] & \unsure{G^{\sharp}|G^{\natural}}
\end{tikzcd}}
\quad\text{i.e.}\quad
\diagram{\begin{tikzcd}[column sep=large,column sep=9pt,row sep=14pt,execute at begin picture={\fontsize{10}{12}\selectfont}]
\unsure{C|G}\arrow[r] & \unsure{G^{\sharp}(S_0)|G^{\natural}(S_0)}
\end{tikzcd}}
\]
satisfaisant les deux conditions suivantes :

a) Commutativité du diagramme
\[
\tag{41}
\diagram{\begin{tikzcd}[column sep=large,row sep=large]
M_{S_0}\arrow[r]\arrow[d] & \unsure{C^+_{S_0}|G^+_{S_0}}\arrow[ddr,"u_{\unsure{C|G}}"] & \\
\underline M_{S_0}\arrow[drr,"u_{\underline M}"'] & & \\
& & G
\end{tikzcd}}
\quad,
\]
et b)
\[
\tag{42}
\begin{gathered}
u_{\underline M}\text{ est compatible avec les actions}\\
\text{de }\unsure{C^+_{S_0}|G^+_{S_0}}\text{ sur }\underline M_{S_0}
\text{ (via son action sur }M_{S_0}\text{)}\\
\text{et sur }G\text{ (via }u_{\unsure{C|G}}\text{).}
\end{gathered}
\]
Dans le cas actuel, on aura $G=\unsure{\underline C|\underline G}$
\[
\tag{43}
u_{\underline M}(\lambda)=\deleted{\unsure{\lambda}}\,i(\mu.\lambda)
\]
où
\[
i:\quad
\diagram{\begin{tikzcd}[column sep=large]
\underline M\arrow[r,hook] & \unsure{\underline C^+|\underline G^+}
\end{tikzcd}}
\]
est l'inclusion, et
\[
\tag{44}
u_{\unsure{C|G}}:\deleted{\text{défini par}}
\inserted{\diagram{\begin{tikzcd}[column sep=large]
\unsure{C^+_{S_0}|G^+_{S_0}}\arrow[r]&G
\end{tikzcd}}}
\qquad\text{défini par (33),}
\]
Il faut donc d'abord vérifier que (33), définit bien un hom
\diagram{\begin{tikzcd}[column sep=small]\unsure{C_{S_0}|G_{S_0}}\arrow[r]&G\end{tikzcd}},
ou ce qui revient au $\overline{\mathrm m}$,
\diagram{\begin{tikzcd}[column sep=small]\unsure{C|G}\arrow[r]&G(S_0)\end{tikzcd}}.
\margin{43}
\margin{43\\\texttt{https://grothendieck.umontpellier.fr}\\Université de Montpellier, UMR CNRS 5149, Case courrier 051, Place Eugène Bataillon\\34095 Montpellier cedex 5 - France}

\scanpage{45}
\mspage{\textcircled{670}}
Pour cela on \deleted{\unsure{note que} $u$ définit un hom \unsure{au} \unsure{d'indice} de $3$
$\unsure{C|G}$}
\inserted{utilise la présentation \deleted{\unsure{de}} de $\unsure{C^+|G^+}$}
\[
\deleted{\begin{gathered}
\unsure{C_0^+|G_0^+}\\
\unsure{\simeq}\\
\unsure{\operatorname{SL}(2,\mathbb Z)}=
\{\rho,\sigma\mid\rho^6=\sigma^4=1\}
\end{gathered}}
\]
\[
\tag{44}
\unsure{C^+|G^+}=
\left\{\rho,\sigma\ \middle|\ 
\begin{gathered}
\rho^6=\sigma^4=1,\quad\rho^3=\sigma^2\\
[\sigma,\lambda_\infty]=1\quad\text{où}\quad
\lambda_\infty=\deleted{\unsure{\rho}}\rho\sigma\rho^{-1}\sigma\rho
\end{gathered}
\right\},
\]
donc posant
\[
\tag{45}
\left\{\begin{aligned}
\rho'&=\unsure{\zeta\rho|u(\rho)},\quad
\sigma'=\eta_0\omega^{\nu_2}\sigma\qquad\text{cf (33) et}\\
\unsure{\zeta|\xi}&=\alpha_0-\rho(\alpha_0)+\nu_3\lambda_1,\quad
\unsure{\eta|\eta_0}=\beta_0-\sigma(\beta_0)
\end{aligned}\right.
\]
il faut vérifier les relations
\[
\tag{46}
\rho'{}^6=\sigma'{}^4=[\sigma',\lambda'_\infty]=1,\quad
\rho'{}^3=\sigma'{}^2
\quad\text{où}\quad
\lambda'_\infty=\deleted{\unsure{\rho'(\lambda'_0)}}
\inserted{\rho'\sigma'\rho'{}^{-1}\sigma'\rho'}.
\]
On commencera par vérifier \deleted{\unsure{que|plus}}
\[
\tag{47}\rho'{}^3=\sigma'{}^2=\omega
\]
ce qui implique \deleted{\unsure{que}}
\[
\tag{48}
\rho'{}^6=\sigma'{}^4=1,\quad\text{donc}\quad
\exists\ \deleted{\text{unique}}
\diagram{\begin{tikzcd}[column sep=large]
\unsure{C_0^+|G_0^+}\arrow[r,"u_{\unsure{C_0|G_0}}"] & G(S_0)\\
\rho\arrow[r,mapsto]&\rho',\quad\sigma\mapsto\sigma'
\end{tikzcd}}
\]
Ensuite, \deleted{posant} notant qu'on aura
\[
\lambda'_\infty=u_{\unsure{C_0|G_0}}(\lambda_\infty),
\]
il suffira de \deleted{\unsure{cela}} vérifier
\[
\tag{49}\lambda'_\infty=\mu\lambda_\infty
\]
il sera évident que $\lambda'_\infty$ commute à
\[
\sigma'=u(\sigma)=\eta_0\sigma\qquad\deleted{\text{\unsure{pui}}}
\]
puisque $\sigma$ et $\eta_0\in M$ y commutent. Donc pour construire
$u_{\unsure{C|G}}$, les vérifications essentielles sont (47) et (49)
--- d'ailleurs une fois (47) prouvé, les calculs précédents aboutissant à (38)
sont justifiés (\unsure{au moins} là où ils ont un sens, i.e.,
pour $u_{\unsure{C|G}}$ et $\lambda\in M$),
\margin{44}
\margin{44\\\texttt{https://grothendieck.umontpellier.fr}\\Université de Montpellier, UMR CNRS 5149, Case courrier 051, Place Eugène Bataillon\\34095 Montpellier cedex 5 - France}

\scanpage{46}
donc (49) sera bien vérifié.

On \deleted{aura} a
\[
\begin{aligned}
&\deleted{\unsure{u(\rho)=}}\\
\rho'{}^3&=\bigl(\unsure{\zeta|\xi}\rho(\unsure{\zeta|\xi})\rho^2(\unsure{\zeta|\xi})\bigr)\rho^3
=\unsure{\zeta|\xi}\rho(\unsure{\zeta|\xi})\rho^2(\unsure{\zeta|\xi})\,\omega\\
\sigma'{}^2&=(\eta_0\sigma(\eta_0))\sigma^2=(\eta_0\sigma(\eta_0))\omega
\end{aligned}
\]
donc \deleted{et} les formules (47) se réduisent à :
\[
\tag{50}
\unsure{\zeta|\xi}\rho(\unsure{\zeta|\xi})\rho^2(\unsure{\zeta|\xi})=1,
\qquad\eta_0\sigma(\eta_0)=1
\]
soit, additivement
\[
\unsure{\zeta|\xi}+\rho(\unsure{\zeta|\xi})+\rho^2(\unsure{\zeta|\xi})=0,
\qquad\eta_0+\sigma(\eta_0)=0.
\]
dont la vérification est triviale, compte tenu que
\[
\rho_{\underline M}^3=\sigma_{\underline M}^2=\mathrm{id}_{\underline M},
\qquad(1+\rho_{\underline M}+\rho_{\underline M}^2)
=\deleted{\unsure{\mathrm{id}_{\underline M}}}0_{\underline M}.
\]
On a donc bien construit un hom
\deleted{\unsure{au}}
\[
u_{\unsure{C_0^+|G_0^+}}:\quad
\diagram{\begin{tikzcd}[column sep=large]
\unsure{C_0^+|G_0^+}\arrow[r]&\unsure{\underline C_0^+(S_0)|\underline G_0^+(S_0)}
\end{tikzcd}},
\]
et il reste à vérifier (41) et (42).

Pour (41), nous l'avons déjà montré au \inserted{niveau de l'existence d'un $u_{\unsure{C_0|G_0}}$}
--- elle revient à prouver
\[
\bigl(\quad u_{\unsure{C_0|G_0}}(\lambda_0)=\deleted{\unsure{\lambda_0}}\mu\lambda_0,
\quad u_{\unsure{C_0|G_0}}(\lambda_1)=\mu\lambda_1
\]
Je vais quand $\overline{\mathrm m}$ refaire la vérification, pour être sûr.
\[
\begin{aligned}
\varepsilon'_0&=u_{\unsure{C_0|G_0}}(\underbrace{\varepsilon_0}_{\sigma\rho})
=u_{\unsure{C_0|G_0}}(\sigma)u_{\unsure{C_0|G_0}}(\rho)
=\eta_0\omega^{\nu_2}\deleted{\unsure{\zeta|\xi}}\sigma\unsure{\zeta|\xi}\rho
=\eta_0\omega^{\nu_2}\sigma(\unsure{\zeta|\xi})\sigma\rho\\
&=\eta_0\sigma(\unsure{\zeta|\xi})\omega^{\nu_2}\varepsilon_0
\end{aligned}
\]
\margin{45}
\margin{45\\\texttt{https://grothendieck.umontpellier.fr}\\Université de Montpellier, UMR CNRS 5149, Case courrier 051, Place Eugène Bataillon\\34095 Montpellier cedex 5 - France}

\scanpage{47}
\mspage{\unsure{671}}

Je dis que
\[
\begin{aligned}
\varepsilon'_0&=\varepsilon_0^{(\mu)}\quad\text{i.e.}\quad
\eta_0\sigma(\unsure{\zeta|\xi})\omega^{(\nu_2)}
 =\underbrace{\varepsilon_0^{(\mu-1)}}_{l_0^\nu}
\quad\text{où }\mu-1=2\nu\deleted{\unsure{+1|=1}},\quad\text{soit}\\
\eta_0\sigma(\unsure{\zeta|\xi})\omega^{\nu_2}&=l_0^{(\nu)}
\quad\text{où }\nu=2\nu'\unsure{-|+}\nu_2\quad\text{i.e.}\\
\eta_0\sigma(\unsure{\zeta|\xi})
 &=\underbrace{l_0^{(2\nu')}}_{\lambda_0^{2\nu'}}
 \underbrace{\unsure{(\omega l_0)^{(-\nu_2)}|(\omega l_0^{(-\nu_2)})}}_{\lambda_0^{-\nu_2}}
 =\underbrace{\lambda_0^{(2\nu'\unsure{-|+}\nu_2)}}_{\lambda_0^\nu}
 \quad\text{i.e.}\\
\eta_0\sigma(\unsure{\zeta|\xi})&=\lambda_0^{(\nu)}
\end{aligned}
\]
soit, en appliquant $\sigma$ et notant que
$\sigma(\eta_0)=\eta_0^{-1}$, $\sigma^2(\unsure{\zeta|\xi})=\unsure{\zeta|\xi}$,
$\sigma(\lambda_0)=\lambda_1$,
\[
\eta_0^{-1}\unsure{\zeta|\xi}=\lambda_1(\nu)
\quad\text{soit enfin }\deleted{\unsure{\zeta_0=|\xi_0=}}
\unsure{\zeta|\xi}=\eta_0\unsure{\lambda_1(\nu)|\lambda_1^{(\nu)}}
\quad\text{ou}
\]
\[
\unsure{\zeta|\xi}=\eta_0+\nu\lambda_1\quad\text{(écrit additivement)}
\]
ce qui est vrai par hypothèse (cf (31)).

On aura ensuite
(\deleted{\unsure{pour $u=u_{\unsure{C_0|G_0}}$|pour $u_{\unsure{C_0|G_0}}$}}
\inserted{écrivant $u$ au lieu de $u_{\unsure{C_0|G_0}}$ pour simplifier})
\[
u_{\deleted{\unsure{C_0|G_0}}}(\lambda_1)
 =u\rho(\lambda_0)
 =(\unsure{\zeta|\xi}\rho u)(\lambda_0)
 =\unsure{\zeta|\xi}\rho(u(\lambda_0))
 =\unsure{\zeta|\xi}\rho(\lambda_0(\mu))
 =\lambda_1(\mu)
\]
OK.

Reste à vérifier (42), ce qui \deleted{\unsure{signifie}}
\inserted{revient} à la \unsure{commutativité} de
$u_{\underline M}=\mu\,\mathrm{id}_{\underline M}$ à
$(\rho_{\underline M},\rho'_{\underline M})$,
$(\sigma_{\underline M},\sigma'_{\underline M})$
--- \unsure{ce signifie} aussi que
$\rho_{\underline M}$ et $\rho'_{\underline M}$ d'une part,
$\sigma_{\underline M}$ et $\sigma'_{\underline M}$ d'autre part,
\deleted{\illegible{2}}\inserted{sont égaux, ce}
qui est également évident.

Ainsi, le groupe $\widetilde Z_{\rho,\sigma}$ est bel et
bien construit --- du moins, en tant que
schéma sur $S_0=\mathrm{Spec}\,\mathbb Z$, avec des morphismes
de schémas
\[
\tag{51}
\diagram{$
\begin{tikzcd}[column sep=3em,row sep=.25em,nodes in empty cells]
\widetilde{\underline Z}_{\rho,\sigma}\arrow[r,"\nu'"] & G_a\\
{}\arrow[r,"\deleted{\unsure{\nu|\varepsilon}}"] & \deleted{\unsure{\{\pm1\}_{S_0}^{2}|\mu_{S_0}^{2}}}\\
{}\arrow[r,"\nu_2"] & \{0,1\}_{S_0}\\
{}\arrow[r,"\nu_3"] & \{0,1\}_{S_0}
\end{tikzcd}
\quad\Bigg|\quad
\begin{gathered}
\text{d'où }\mu=4\nu'+2+\nu_2:\quad
\begin{tikzcd}[column sep=1.5em,row sep=1em]
\widetilde Z_{\rho,\sigma}\arrow[r]\arrow[rd] & G_a\\
 & G_m\arrow[u,hook]
\end{tikzcd}\\
\varepsilon_2,\varepsilon_3:\quad
\begin{tikzcd}[column sep=3em]
\widetilde Z_{\rho,\sigma}\arrow[r,shift left=.3ex]\arrow[r,shift right=.3ex]
 & \mu_{S_0}=\{\pm1\}_{S_0}
\end{tikzcd}
\end{gathered}
$}
\]

\margin{46}
\margin{46\\\texttt{https://grothendieck.umontpellier.fr}\\Université de Montpellier, UMR CNRS 5149, Case courrier 051, Place Eugène Bataillon\\34095 Montpellier cedex 5 - France}

\scanpage{48}

et un morphisme de \deleted{\unsure{foncteurs}} \inserted{préfaisceaux} sur $S_0$
\[
(52)\qquad
\diagram{\begin{tikzcd}[row sep=small]
\widetilde{\underline Z}_{\rho,\sigma} \arrow[r] & \underline{\operatorname{Aut}}_{\mathrm{gr}}\bigl(\unsure{\underline C^{+}|\underline G^{+}}\bigr),\\
\underline u \arrow[r,mapsto] & u
\end{tikzcd}}
\]
donnant lieu à un morphisme
\[
(53)\qquad
\diagram{\begin{tikzcd}[row sep=small]
\widetilde{\underline Z}_{\rho,\sigma} \arrow[r] & \underline{\operatorname{Aut}}_{\mathrm{gr}}\bigl(\unsure{\underline C^{+}|\underline G^{+}}\bigr)\times\mathbf G_{\mathrm m}\times\mu_{S_0}\times\mu_{S_0}\\
\underline u \arrow[r,mapsto] & u,\mu(\underline u),\varepsilon_2(\underline u),\varepsilon_3(\underline u)
\end{tikzcd}}
\]
--- des calculs déjà faits trente six fois dans des contextes divers doivent montrer que ce morphisme (qui est \uline{mono.}\ \unsure{?}\ en car. $\ne 2,3$) est un morphisme de \uline{groupes}, pour une structure de \uline{\unsure{groupe|groupes}} convenable sur $\widetilde{\underline Z}_{\rho,\sigma}$, uniquement déterminée (car elle l'est en car. $\ne 2,3$, et on conclut par raison de densité schématique). On aura \deleted{\unsure{ce n'est pas}}
\[
(54)\qquad
\diagram{\begin{tikzcd}
1\arrow[r]&\widetilde{\underline Z}_{\rho,\sigma}^{+0}\arrow[r]&\widetilde{\underline Z}_{\rho,\sigma}\arrow[r]&\mathbf G_{\mathrm m}\times\mu_{S_0}\times\mu_{S_0}\arrow[r]&1
\end{tikzcd}}
\]
\margin{un peu \unsure{cavalier}, car
\diagram{\begin{tikzcd}
\underline Z_{\rho,\sigma}\arrow[r]&\mathbf G_{\mathrm m}\times\mu_{S_0}\times\mu_{S_0}
\end{tikzcd}}
n'est pas un épi, du fait que
\diagram{\begin{tikzcd}[row sep=small]
\underline M\times\underline M\arrow[r]&\underline M\\
(\alpha_0,\beta_0)\arrow[r,mapsto]&(1-\rho_{\underline M})\alpha_0\unsure{-|+}(1-\sigma_{\underline M})\beta_0
\end{tikzcd}}
n'est pas épi --- \unsure{où il} y a \unsure{canular} en car. 3 --- cf plus bas\ldots}
en explicitant
\[
(55)\qquad
\begin{gathered}
\widetilde{\underline Z}_{\rho,\sigma}^{+0}\underset{\inserted{\text{sous-schéma en groupes}}}{\subset}\underline M\times\underline M\times\mathbf G_{\mathrm a}\\
\Vert\\
\left\{(\alpha_0,\beta_0,\nu')\,\middle|\,
\begin{gathered}
\deleted{\unsure{\alpha_0}}4\nu'=0\\
(1+\rho_{\underline M})\alpha_0=(1-\sigma_{\underline M})\beta_0+2\nu'\lambda_1
\end{gathered}\right\}
\end{gathered}
\]
On aura donc
\[
(56)\qquad
\diagram{\begin{tikzcd}
1\arrow[r]&\widetilde{\underline Z}_{\rho,\sigma}^{+00}\arrow[r]&\widetilde{\underline Z}_{\rho,\sigma}^{+0}\arrow[r]&\underline{{}_4\mathbf G_{\mathrm a}}\arrow[r]&1
\end{tikzcd}}
\]
\margin{47}
\margin{47\\\texttt{https://grothendieck.umontpellier.fr}\\Université de Montpellier, UMR CNRS 5149, Case courrier 051, Place Eugène Bataillon\\34095 Montpellier cedex 5 - France}

\scanpage{49}
\mspage{672}

où
\[
(57)\qquad
\begin{gathered}
\widetilde{\underline Z}_{\rho,\sigma}^{+00}\overset{\mathrm{dfn}}{=}
\left\{\underline u\in\Gamma\widetilde{\underline Z}_{\rho,\sigma}\,\middle|\,
\nu'(\underline u)=\nu_2(\underline u)=\nu_3(\unsure{u|\underline u})=0\right\}\\
\simeq\\
\left\{(\alpha_0,\beta_0)\in\Gamma(\underline M\times\underline M)\,\middle|\,
(1-\rho_{\underline M})\alpha_0=(1-\sigma_{\underline M})\beta_0\right\}.
\end{gathered}
\]
\margin{en car. 3 sa dimension relative \unsure{saute} de 2 à 3 !}
schéma en \uline{groupes} $\underline Z_{\rho,\sigma}^{+00}$ est lisse \inserted{en car. $\ne 3$} de dim relative 2 --- c'est $\overline{\mathrm m}$ de façon canonique un schéma en Modules, on aura
\[
(58)\qquad\widetilde{\underline Z}_{\rho,\sigma}^{+00}\simeq\unsure{\mathbf V|\mathbb V}(\underline P)
\]
où $P$ est le $\mathbb Z$-module libre de rang 2 défini \deleted{\illegible{1}} par la suite exacte
\[
(59)\qquad
\diagram{\begin{tikzcd}[row sep=small]
0\arrow[r]&\deleted{\unsure{M}}\inserted{P}\arrow[r]&M\times M\arrow[r]&M\arrow[r]&0\\
&&(\alpha_0,\beta_0)\arrow[r,mapsto]&(1-\rho_{\underline M})\alpha_0-(1-\sigma_{\underline M})\beta_0&
\end{tikzcd}}
\]
L'hom $M\times M\longrightarrow M$ s'explicite, en posant
\[
(60)\qquad
\left\{\begin{aligned}
\alpha_0&=a_0\lambda_0+\unsure{a_1}\lambda_1
&\text{d'où}\quad\alpha_0-\rho\alpha_0&=(a_0+\unsure{a_1})\lambda_0+(2\unsure{a_1|b_1}-a_0)\lambda_1\\
\beta_0&=b_0\lambda_0+b_1\lambda_1
&\text{d'où}\quad\beta_0-\sigma\beta_0&=(b_0-b_1)\lambda_0+(b_1-b_0)\lambda_1
\end{aligned}\right.
\]
comme
\[
(61)\qquad
\diagram{\begin{tikzcd}
(\alpha_0,\beta_0)\arrow[r,mapsto]&
\underbrace{(a_0+a_1+(b_1\unsure{-|+}b_0))}_{c_0}\lambda_0+
\underbrace{(2a_1-a_0-(b_1-b_0))}_{c_1}\lambda_1
\end{tikzcd}}
\]
\deleted{\unsure{On}} C'est un hom surjectif \inserted{sauf en car. 3}, \deleted{pour $(c_0,c_1)$ donnés,}
\deleted{\unsure{on fera}\ $\unsure{a_1|c_1}$\ \illegible{1}}
\inserted{--- les éléments de l'image sont les}
\[
(62)\qquad
\left\{c_0\lambda_0+c_1\lambda_1\,\middle|\,
\underbrace{c_0+c_1}_{=3a_1}\equiv0\quad(3)\right\}
\]
\uline{Remarque} Si on prend \deleted{\unsure{un}}\inserted{un} $u$ défini sur $\widehat{\mathbb Z}$ provenant d'un élément de $M[0]$, alors
\margin{48}
\margin{48\\\texttt{https://grothendieck.umontpellier.fr}\\Université de Montpellier, UMR CNRS 5149, Case courrier 051, Place Eugène Bataillon\\34095 Montpellier cedex 5 - France}

\scanpage{50}

on aura \deleted{que}
\[
\mu\in\widehat{\mathbb Z}^{*},\quad
\mu\equiv\deleted{\unsure{\varepsilon_3|\varepsilon_2}}(-1)^{\nu_3}(=1-2\nu_3)\quad\mathrm{mod}\ (3)
\]
et la condition sur $(\mu,\varepsilon_2,\varepsilon_3)$ pour être dans l'image de
\[
\diagram{\begin{tikzcd}
\widetilde{\underline Z}_{\rho,\sigma}\arrow[r]&\mathbf G_{\mathrm m}\times\mu_{S_0}\times\mu_{S_0}
\end{tikzcd}},
\]
qui s'écrit en général (en vertu de (62) \unsure{\&} et de (32))
\[
2\nu'-\nu_2+\nu_3\equiv0\quad(3)
\]
devient, en multipliant par 2
\[
\underbrace{4\nu'-2\nu_2}_{\mu-1}+2\nu_3\equiv0\quad(3)
\]
i.e
\[
\mu\equiv1-2\nu_3\quad(3),
\]
et elle est automatiquement satisfaite.

\deleted{
C'est pourquoi j'ai envie de modifier la définition de \deleted{$\unsure{Z}$} $\widetilde{\underline Z}_{\rho,\sigma}$, en définissant $\underline Z_{\rho,\sigma}$ comme
\[
\deleted{(63)\quad\widetilde{\underline Z}_{\rho,\sigma}=}
\left\{\begin{gathered}
\deleted{\unsure{\Gamma|\mu}}(\inserted{\mu,}\nu',\nu'',\nu_2,\nu_3,\alpha_0,\beta_0)
\in\mathbf G_{\mathrm m}\times\mathbf G_{\mathrm a}\times\unsure{\mathbf G_{\mathrm a}}\times\{0,1\}_{S_0}\times\{0,1\}_{S_0}\times\underline M\times\underline M\\
\left|\begin{gathered}
\mu=4\nu'\unsure{-|+}2\unsure{\nu_2}+1(=6\nu''-2\nu_3+1),\quad 2\nu'+\nu_3-\nu_2=3\nu'',\\
(\unsure{1}-\rho_{\underline M})(\alpha_0)-(1-\sigma_{\underline M})(\beta_0)
=\underbrace{(\deleted{\unsure{2\nu'-\nu_2+\nu_3}}\inserted{3\nu''})}_{c}\lambda_1
\end{gathered}\right.
\end{gathered}\right\}
\]
}
\margin{en \unsure{exprimant|exigeant} $\unsure{\mu|\nu}$\ \unsure{s'écrivent}
\[
\begin{aligned}
\mu&\equiv1-2\nu_2\quad(4)\\
\mu&\equiv1-2\nu_3\quad(6)
\end{aligned}
\unsure{\left[\vphantom{\begin{matrix}\mu\\\mu\end{matrix}}\right.}
\]
\deleted{l'équation des lacets}
\[
\begin{aligned}
\mu&=1\unsure{-|+}2\nu_2+4\nu'&&\text{\unsure{avec}}\\
\mu&=1-2\nu_3+6\nu''\\
-\nu_2+2\nu'&=-\nu_3+3\nu''&&\text{i.e.}\\
2\nu'+\nu_3-\nu_2&=3\nu''
\end{aligned}
\]}
\deleted{
\deleted{On aura ici
\[
2c=6\nu'',\quad\text{et on \unsure{va|a} en fait}
\]
prendre $c=3\nu''$, et.}
On aura
\[
(63)\qquad
\begin{gathered}
\widetilde{\underline Z}_{\rho,\sigma}\subset\mathbf G_{\mathrm a}\times\mathbf G_{\mathrm a}\times\{0,1\}_{S_0}\times\{0,1\}_{S_0}\times\underline M\times\underline M\\
\Vert\\
\left\{(\nu',\nu'',\nu_2,\nu_3,\alpha_0,\beta_0)\,\middle|\,
\begin{aligned}
\text{a)}\quad &2\nu'\deleted{+\unsure{\nu_3}}-\nu_2=3\nu''-\nu_3\\
&\deleted{(\text{d'où }4\nu'\unsure{-|+}2\nu_2+1=6\nu''-2\nu_3+1,\ \text{soit }\mu)}\\
&\inserted{\text{soit }\nu,\ \text{d'où }\mu=2\nu+1}\\
\text{b)}\quad &\mu\in\Gamma\mathbf G_{\mathrm m}\\
\text{c)}\quad &\text{équation des \unsure{lacets|lacet}}\\
&(1-\rho_{\underline M})(\alpha_0)-(1-\sigma_{\underline M})(\beta_0)=3\nu''\unsure{\underline{\lambda_1}|\lambda_1}
\end{aligned}\right\}
\end{gathered}
\]
}
\margin{49}
\margin{49\\\texttt{https://grothendieck.umontpellier.fr}\\Université de Montpellier, UMR CNRS 5149, Case courrier 051, Place Eugène Bataillon\\34095 Montpellier cedex 5 - France}

\scanpage{51}
\mspage{673}

3. Les groupes \deleted{$\unsure{\widetilde Z_{\rho,\sigma}}$ \unsure{révisés} \illegible{2} $\mathbb O_N$} $\mathbb O_N^{*}$.

On va introduire une construction auxiliaire,
en définissant un schéma en groupes qui
correspond à la donnée de quantités
$\mu\in\Gamma\mathbf G_{\mathrm m}$, $\nu',\nu''\in\Gamma\mathbf G_{\mathrm a}$,
$\nu_2,\nu_3\in\Gamma\{0,1\}_{S_0}$, satisfai-
sant
\[
(63)\qquad
\mu=4\nu'-2\nu_2+1=6\nu''-2\nu_3+1,
\qquad 2\nu'-\nu_2=3\nu''-\nu_3.
\]

De façon générale, soit $N$\inserted{$\in\mathbb N^{*}$} un entier \inserted{naturel} fixé
(le cas qui nous intéresse pour la suite est
$N=12$), et considérons le schéma \inserted{en groupes, $\mathbb O_N$}
sur $\mathbb Z$ associé à l'extension de groupes discrets
\deleted{\[(64)\qquad \mathbb O_N\]}
\[
(64)\qquad
\diagram{\begin{tikzcd}[column sep=large]
1\arrow[r]&\mathbb Z\arrow[r,"N"]&\mathbb Z\arrow[r]&\mathbb Z/N\mathbb Z\arrow[r]&1
\end{tikzcd}}
\]
qui s'insère dans une suite exacte
\[
(65)\qquad
\diagram{\begin{tikzcd}[column sep=large]
1\arrow[r]&\unsure{\mathbf E^{1}|\mathbb E^{1}}\arrow[r,"i"]&\mathbb O_N\arrow[r]&(\mathbb Z/N\mathbb Z)_{S_0}\arrow[r]&1.
\end{tikzcd}}
\]

Rappelons la \deleted{\unsure{défi}} construction simple
\[
(66)\qquad
\mathbb O_N
=(\mathbb Z)_{S_0}\mathop{\wedge}_{(N\mathbb Z)}\underline{N\mathbb Z}.
\]
\deleted{Mais} Je dis qu'il y a \unsure{dans|sur} $\mathbb O_N$ une
structure d'Anneaux unique, dont l'
addition soit la structure de groupe \inserted{abélien} déjà
envisagée, et où la multiplication est
obtenue ainsi : Une section de $\mathbb O_N$ est
\deleted{locale} de la forme \deleted{$\unsure{k|x}$} $n+N[\lambda]$, avec
les règles de calcul
\[
(67)\qquad
\left\{\begin{aligned}
\text{a)}\quad &(n+N[\lambda])+(n'+N[\lambda'])=(n+n')+N[\lambda+\lambda']\\
\text{b)}\quad &n+N[\lambda]=n'+N[\lambda']\Longleftrightarrow
\deleted{\unsure{\xi|n}}\ n+N\lambda=n'+N\lambda'\text{ dans }\Gamma\mathbf G_{\mathrm a}.
\end{aligned}\right.
\]

\margin{50}
\margin{50\\\texttt{https://grothendieck.umontpellier.fr}\\Université de Montpellier, UMR CNRS 5149, Case courrier 051, Place Eugène Bataillon\\34095 Montpellier cedex 5 - France}

\scanpage{52}

On multiplie alors les sections par
\[
(68)\qquad
(\unsure{n}+N[\lambda])(n'+N[\lambda'])
=nn'+N[n\lambda'+n'\lambda+\lambda\lambda'],
\]
on vérifie que c'est \deleted{\unsure{biad}} compatible aux
relations (67 b), que c'est biadditif,
commutatif, associatif, que $1+N[0]$ est
unité\ldots{} Autre façon de \deleted{\unsure{la|le} \unsure{con|en}} voir la
structure d'anneaux. \deleted{On} Soit plus
généralement \deleted{$\mathbb O_N=\mathbb Z_{S_0}$}
$A$ une \deleted{$\unsure{\mathbb Z|z}$} $\mathbb Z$-algèbre
\deleted{\unsure{com}} associative \inserted{\unsure{unifère|unitaire}} (disons)
\deleted{\unsure{principal}} \inserted{munie d'un \uline{idéal}}
\margin{\inserted{bilatère}} $J$ \inserted{qui soit}
un $\mathbb Z$-module libre de t.f., considérons
\deleted{\unsure{un} idéal}
\deleted{l'idéal $NA$} et la structure d'
extension
\[
(69)\qquad
\diagram{\begin{tikzcd}[column sep=large]
0\arrow[r]&\deleted{\unsure{NA}}\inserted{J}\arrow[r]&A\arrow[r]&\overbrace{A/\deleted{\unsure{NA|N}}\inserted{J}}^{B}\arrow[r]&0
\end{tikzcd}}
\]
d'où un schéma en groupes commu-
tatif $\underline A_{[N]}$, qui s'obtient comme
quotient de
\deleted{$(A)_{\unsure{S|S_0}}\cdot\underline{\unsure{NB|NA}}\;\unsure{\xi\to|p}$}
\inserted{$A_{S_0}\cdot\underline J$ produit}
semi-direct)
\[
(70)\qquad
\diagram{\begin{tikzcd}[column sep=large]
0\arrow[r]&\underline J\arrow[r]&A_{S_0}\times\underline J\arrow[r]&\underline A_N\arrow[r]&0\\
&\lambda\arrow[r,mapsto]&(\lambda,-\lambda)&&
\end{tikzcd}}
\]
et on trouve que $A_{S_0}\cdot\underline J$ a une structure
\uline{d'Anneaux} produit $\frac12$-direct, et que

\margin{51}
\margin{51\\\texttt{https://grothendieck.umontpellier.fr}\\Université de Montpellier, UMR CNRS 5149, Case courrier 051, Place Eugène Bataillon\\34095 Montpellier cedex 5 - France}

\scanpage{53}
\mspage{674}

l'image de $\underline J$ y est un idéal ...

On trouve bien une structure d'anneau sur $\underline A_{[N]}$, et une suite exacte \deleted{\unsure{de}}
\[
\tag{71}
\diagram{\begin{tikzcd}[column sep=large]
1 \arrow[r,no head] & \underline J \arrow[r,no head] & \underline A_{[N]} \arrow[r] & B_{S_0} \arrow[r] & 1
\end{tikzcd}}
\]
où $\underline A_{[N]}\mathbin{-}B_{S_0}$ est un \unsure{hom.} de \uline{schémas en anneaux}. Passant aux groupes \deleted{\unsure{anneaux}} multiplicatifs, on trouve une suite exacte
\[
\tag{72}
\diagram{\begin{tikzcd}[column sep=large,row sep=small]
1 \arrow[r] & (1+\underline J)^* \arrow[r,no head] & A_{[N]}^* \arrow[r,no head] & (B_{S_0})^* \arrow[r,no head] \arrow[d,equal] & 1 \\
& & & (B^*)_{S_0} &
\end{tikzcd}}
\]
où pour $A_{[N]}^*$, $(B_{S_0})^{\unsure{*|'}}$ et $B^*$, le $*$ désigne le \unsure{passage} au \unsure{s-groupe} multiplicatif et où $(1+\underline J)^*$ a le sens analogue
\[
\tag{73}
\diagram{\begin{tikzcd}[row sep=large]
(1+\underline J)^*\unsure{\subset|\subseteq}1+\underline J \arrow[d,draw=none,"="{description,rotate=90}] \\
\{(1+\lambda)\mid 1+\lambda\ \text{inv. dans }\underline A\} \arrow[d,draw=none,"\simeq"{description,rotate=90}] \\
\{(\lambda,\lambda')\in\underline J\times\underline J\mid\lambda+\lambda'+\lambda\lambda'=0\ \text{dans }\underline J\}.
\end{tikzcd}}
\]

Ainsi on aura, dans le cas particulier $A=\mathbb Z$, $J=N\mathbb Z$
\[
\tag{74}
\diagram{\begin{tikzcd}[column sep=large]
1 \arrow[r] & (\underline{1+N\mathbb Z})^* \arrow[r] & \mathbb O_N^* \arrow[r,no head] & (\mathbb Z/N\mathbb Z)^*_{S_0} \arrow[r,no head] & 1
\end{tikzcd}}
\]
où le noyau

\margin{52}
\margin{52\\\texttt{https://grothendieck.umontpellier.fr}\\Université de Montpellier, UMR CNRS 5149, Case courrier 051, Place Eugène Bataillon\\34095 Montpellier cedex 5 - France}

\scanpage{54}
\deleted{\mspage{\unsure{674}}}

$(1+N\mathbb Z)^*$ est défini comme
\[
\tag{75}
(1+N\mathbb Z)^*\overset{\mathrm{dfn}}{=}
\{\lambda\in\deleted{\unsure{\mathbf G_a}}\underline{\mathbb Z}=\mathbf G_a\mid 1+N\lambda\ \text{inversible}\}
\]

\deleted{On a un \unsure{hom.} canonique (\unsure{en partie}}

Les $\mathbb O_N$ pour $N$ variable forment un syst. projectif d'Anneaux, les $\mathbb O_N^*$ forment un syst. projectif de schémas en groupes commutatifs.

Je ne vais pas regarder de plus près la structure de ces systèmes projectifs, ni examiner \unsure{les} limites projectives --- \unsure{questions} qui se posent plus \unsure{générales|généralement} quand on part d'un $A$ comme ci-dessus. On va expliciter l'\unsure{hom.} canonique
\[
\tag{76}
\diagram{\begin{tikzcd}[column sep=large]
\mathbb O_N^* \arrow[r,no head] & \mathbb O_1^*=\mathbb O^*=\mathbf G_m
\end{tikzcd}}
\]
\deleted{$\unsure{\varphi}$}

\deleted{dont le noyau est}
qui en général n'est ni épi, ni mono. Son noyau est \deleted{formé des} \deleted{$\unsure{\varphi}\ \unsure{\lambda}$} isomorphe au \unsure{ssgt.|syst.} des
\unsure{$+$}
\[
\tag{77}
\{1+N[\lambda]\mid 1+N\lambda=1\ \text{i.e.}\ N\lambda=0\}
\]
qui en tant que schéma (\unsure{sans} structure de groupe) est isomorphe à ${}_N\mathbf G_a$,

\margin{53}
\margin{53\\\texttt{https://grothendieck.umontpellier.fr}\\Université de Montpellier, UMR CNRS 5149, Case courrier 051, Place Eugène Bataillon\\34095 Montpellier cedex 5 - France}

\scanpage{55}
\mspage{675}

qui est un schéma en groupes (non plat si $N\ne1$) qui est nul sur $\operatorname{Spec}\mathbb Z[1/N]$, et \deleted{est \unsure{dim.|de dim.} relative $1$} les fibres en les $p\mid N$ sont $\simeq\mathbf G_a$.

Je me rends compte maintenant que c'est une maladresse de sortir \unsure{du} \unsure{groupe} $\mathbb O_N^*$, qui est \unsure{parfait} en son genre, pour se rapporter à $\mathbf G_m$.

Je vais en particulier examiner
\[
\tag{78}
\diagram{\begin{tikzcd}[column sep=large]
1 \arrow[r,no head] & (1+12\mathbb Z)^* \arrow[r,no head] & \mathbb O_{12}^* \arrow[r] & (\mathbb Z/12\mathbb Z)^*_{S_0} \arrow[r,no head] & 1.
\end{tikzcd}}
\]

Notons que \deleted{$\mathbb Z/12\mathbb Z$}
\[
(\mathbb Z/12\mathbb Z)^*\simeq(\mathbb Z/4\mathbb Z)^*\times(\mathbb Z/6\mathbb Z)^*\simeq\mu\times\mu.
\]

Donc (78) prend la forme
\deleted{et on trouve une section ensembliste \unsure{le|la}\ \illegible{1}\ \unsure{de|du}\ \unsure{groupe|noyau}\ \unsure{sur}\ de l'extension $\unsure{\beta}$}
\[
\tag{79}
\diagram{\begin{tikzcd}[column sep=large]
1 \arrow[r,no head] & (1+12\mathbb Z)^* \arrow[r,no head] & \mathbb O_{12}^* \arrow[r,no head,"\varepsilon_2\times\varepsilon_3"] & \mu\times\mu \arrow[r,no head] & 1
\end{tikzcd}}
\]

Les quantités $\nu_{\unsure{2|3}}\ \nu_3$ sont également \inserted{des morphismes} définies sur $\mathbb O_{12}^*$
\margin{\deleted{\unsure{ici}}}
\[
\tag{80}
\diagram{\begin{tikzcd}[column sep=large]
\mathbb O_{12}^* \arrow[r,"(\nu_2,\nu_3)"] & \{0,1\}_{S_0}\times\{0,1\}_{S_0}
\end{tikzcd}}
\]
ce sont les composés
\[
\tag{81}
\diagram{$\begin{tikzcd}[ampersand replacement=\&,column sep=large,row sep=small]
\mathbb O_{12}^* \arrow[r,"(\varepsilon_2,\varepsilon_3)"] \& \mu\times\mu \arrow[r,"\sim"] \& \{0,1\}_{S_0}\times\{0,1\}_{S_0} \\
\& \& \text{déduit de}\quad\deleted{\unsure{(+1)\mapsto(0)}}\quad\begin{array}{rcl}1&\mapsto&0\\-1&\mapsto&1\end{array}
\end{tikzcd}$}
\]

\margin{54}
\margin{54\\\texttt{https://grothendieck.umontpellier.fr}\\Université de Montpellier, UMR CNRS 5149, Case courrier 051, Place Eugène Bataillon\\34095 Montpellier cedex 5 - France}

\scanpage{56}

\deleted{\unsure{que}}

et ce sont des hom de groupes (pour l'addition mod $2$ sur $\{0,1\}_{S_0}$.

\inserted{On a aussi $\varepsilon_2=(-\unsure{1})^{\nu_2}=1-2\nu_2,\quad\varepsilon_3=(-\unsure{1})^{\nu_3}=1-2\nu_3$}

Les ``quatre cas de \uline{signes}'' dans $(\mathbb Z/12\mathbb Z)^*$ sont donnés dans le tableau
\[
\tag{82}
\diagram{$\begin{array}{r|c|c}
{}_{\varepsilon_3}\!\diagdown^{\varepsilon_2}
&\begin{array}{c}+1\\\text{i.e. }\nu_2=0\end{array}
&\begin{array}{c}-1\\\text{i.e. }\nu_2=1\end{array}\\ \hline
\begin{array}{r}+1\\\text{i.e. }\nu_3=0\end{array}&1&+5\\ \hline
\begin{array}{r}-1\\\text{i.e. }\nu_3=1\end{array}&-5&-1\\ \hline
\end{array}$}
\]

donc on trouve une section ensembliste de (79) en prenant les sections de $\underline{\mathbb O}_{12}^*$ de la forme
\[
\tag{83}
\underline{\mu}_0=(-5)^{\nu_2}(+5)^{\nu_3}
\left\{\begin{array}{ll}
\equiv(-1)^{\nu_3}&(6)\\
\equiv(-1)^{\nu_2}&(4)
\end{array}\right.
\]

Donc une section générale de $\mathbb O_{12}^*$ s'écrit
\[
\tag{84}
\underline{\mu}=\underbrace{(-5)^{\nu_2}(+5)^{\nu_3}}_{\mu_0}+\unsure{12}[\unsure{\deleted{\lambda_0}\inserted{\nu_0}|\nu_0}]
\]
où
\[
\left\{\begin{array}{l}
(\nu_2,\nu_3)\in\Gamma(\{0,1\}_{S_0}\times\{0,1\}_{S_0})\\
\lambda_0\in\Gamma(\mathbf G_a).
\end{array}\right.
\]

\deleted{On \unsure{raisonne}\ \unsure{par|plus}\ \unsure{mod}\ $4$}
\inserted{Notons qu'on a}
\[
\deleted{\mu=(4+1}
\qquad
5^{\nu_3}=1+4\nu_3
\]
et
\[
(-5)^{\nu_2}=\left[\begin{aligned}
(-1)^{\nu_2}5^{\nu_2}&=(1-2\nu_2)(1+4\nu_2)\\
&=1+2\nu_2-8\nu_2^2
\end{aligned}\right]=1-6\nu_2
\]
\margin{\uline{NB} $\nu_2=\nu_2^2$}

donc
\[
\tag{\unsure{85}}
\begin{aligned}
\underline{\mu}_0&=(1+4\nu_3)(1\unsure{-6}\nu_2)\deleted{\unsure{-8}\nu_2^2)}\\
&\deleted{=1+2\nu_2+\unsure{(-8)}\nu_2^2+4\nu_3(1+2\nu_2-8\nu_2^2)}\\
&\deleted{=1\unsure{-}2\nu_2+4[\unsure{\nu_2}-2\nu_2^2+\nu_3(1}\\
&=1+4\nu_3-6\nu_2-24\nu_2\nu_3
\end{aligned}
\]

\deleted{On \unsure{pourrait} prendre}
\[
\deleted{(\unsure{85'|85})\qquad\mu'_0=1+4\nu_3-6\nu_2}
\]

\margin{55}
\margin{55\\\texttt{https://grothendieck.umontpellier.fr}\\Université de Montpellier, UMR CNRS 5149, Case courrier 051, Place Eugène Bataillon\\34095 Montpellier cedex 5 - France}

\scanpage{57}
\mspage{676}

Si on \unsure{a} ((84) et (85))
\[
\underline{\mu}=\underline{\mu}_0+12[\unsure{\deleted{\lambda_0}\inserted{\nu_0}|\nu_0}]
=1+4\nu_3-6\nu_2+12(\nu_0-2\nu_2\nu_3)
\]
on en conclut (en \unsure{s'intéressant} aux congruences mod $4$)
\[
\tag{86}
\mu=1-2\nu_2+4\underbrace{[\nu_3-\nu_2\unsure{-|+}3\nu_2\nu_3+3\nu_0]}_{\nu'}
\]
et \unsure{par|pour} \unsure{congruence|congruences} mod $6$
\[
\tag{87}
\mu=1-2\nu_3+6\underbrace{[\nu_3-\nu_2-2\nu_2\nu_3+2\nu_0]}_{\nu''},
\]
où le souligné dans $\underline{\mu}$ désigne la section en $\mathbb O_{12}^*$, le non souligné \inserted{$\mu$} \deleted{\unsure{une}} \inserted{la} section \inserted{\unsure{correspondante|correspondant}} de $\mathbf G_m=\mathbb O_1^*$. On aura bien
\[
\begin{aligned}
2\nu'-\nu_2&=2\nu_3-3\nu_2-6\nu_2\nu_3+6\nu_0\qquad\unsure{''}\\
3\nu''-\nu_3&=2\nu_3-3\nu_2-6\nu_2\nu_3+6\nu_0
\end{aligned}
\]

En résumé

\uline{Proposition} Introduisons sur \inserted{le schéma en \unsure{groupe|groupes} lisse} $\mathbb O_{12}^*$ les fonctions \inserted{$\underline{\mu}\mapsto$} $\unsure{\nu_2|\xi_2},\nu_3\unsure{,}$ \deleted{$\unsure{\xi_3|\varepsilon_3}$ \unsure{et}} (à valeurs dans $\{0,1\}_{S_0}\subset\inserted{\mathbf G_a=}\mathbb O_1$), $\nu',\nu''$ \deleted{$\unsure{\mu_3|\nu_3}$} (à valeurs dans $\unsure{\mathbb O_1|\mathbf G_a}$) et $\mu$ (à valeurs dans $\mathbf G_m\inserted{^{\unsure{*}}=\mathbb O_1^*\subset\mathbb O_1}$) \unsure{définies} par les conditions
\[
\tag{88}
\left\{\begin{array}{ll}
\mu={}&\text{\unsure{image} can. de }\underline{\mu}\text{ dans }\mathbb O_1^*\text{ par}\\
&\text{l'hom de Groupes }\mathbb O_N^*\unsure{\longrightarrow|\mathrel{-}}\mathbb O_1^*\\[4pt]
\begin{gathered}
\inserted{(-1)^{\nu_2}}\\[-2pt]
\deleted{\text{\unsure{l'image}}}\\[-2pt]
\deleted{(\unsure{\varepsilon_2|\xi_2})}
\end{gathered}={}&\text{\unsure{image} can. de }\underline{\mu}\text{ dans }(\mathbb Z/4\mathbb Z)^*_{S_0}\\[4pt]
(-1)^{\nu_3}={}&\rule{5cm}{0.4pt}\quad(\mathbb Z/6\mathbb Z)^*_{S_0}
\end{array}\right.
\]
d'où
\[
\tag{89}
\deleted{\unsure{\Bigg\{}}\quad
\left\{\begin{aligned}
\underline{\mu}&=\deleted{\unsure{\mu_0}}(-5)^{\nu_2}(+5)^{\nu_3}+12[\unsure{\nu_0|\lambda_0}]\\
&=(1+4\nu_3)(1-6\nu_2)+12\nu_0\\
&=1+4\nu_3-6\nu_2+12\unsure{[}\nu_0-2\nu_2\nu_3]
\end{aligned}\right.
\]
pour une section $\nu_0$ bien déterminée de $\mathbf G_a=\mathbb O_1$

\margin{56}
\margin{56\\\texttt{https://grothendieck.umontpellier.fr}\\Université de Montpellier, UMR CNRS 5149, Case courrier 051, Place Eugène Bataillon\\34095 Montpellier cedex 5 - France}
